\chapter{Top-Down Synthesis for Global Temporal Tasks}
\label{chap:top-down}

The previous Chapter~\ref{chap:bottom-up} considered a bottom-up coordination framework in which
temporal tasks are assigned locally and coordination emerges through
execution-time dependence among the robots. That perspective is natural when
task coupling is sparse and collaboration is triggered only when local plans
intersect through shared actions, resources, or newly acquired information.
In many applications, however, the mission is specified from the outset at the
level of the whole team. The central problem is then no longer how local tasks
can be coordinated after they are assigned, but how a global temporal
specification should be decomposed, distributed, and refined so that the
resulting robot behaviors remain jointly correct. This top-down setting is
conceptually complementary to the bottom-up perspective of the previous
chapter, but it raises a different technical difficulty. The logical structure
of the team task must be preserved during decomposition and assignment, while
the resulting planning problem must remain computationally manageable for
heterogeneous teams and large workspaces.

This chapter develops a top-down synthesis framework for such global temporal
tasks. Section~\ref{ch04:sec:po_decomp_assignment} begins with
partial-order-based task decomposition and assignment, where the accepted
behavior of the global task automaton is converted into subtask occurrences
together with a compact temporal-dependence structure, which then supports
efficient allocation over the robot team. Section~\ref{ch04:subsec:task}
extends this perspective to hierarchical coordination of large-scale fleets,
where abstraction across multiple planning layers is used to reduce
computational growth while preserving mission-level consistency.
Section~\ref{ch04:sec:online_continual_tasks} then addresses online adaptation
for continual tasks, in which new demands may arise during execution and the
team plan must be revised without losing the global temporal structure. The
case study and summary at the end illustrate how these
components fit together within a unified top-down coordination framework.


%==============================
%==============================
\section{Partial-Order-Based Task Decomposition and Assignment}
\label{sec:po_decomp_assignment}

This section develops the first component of the top-down synthesis framework
for global temporal tasks. The starting point is the pruned task automaton
$\mathcal{B}_{\varphi}^{-}$ associated with the global sc-LTL formula
$\varphi$. The main objective is to transform the accepted automaton behavior
into a collection of executable subtask occurrences together with a compact
description of their temporal dependence. This description should be rich
enough to preserve correctness with respect to the original task, while being
light enough to support efficient assignment over a multi-robot team.

The key idea is to separate logical structure from allocation decisions.
An accepting run of $\mathcal{B}_{\varphi}^{-}$ provides a valid symbolic
execution of the global task, but the induced order is usually too restrictive.
If this linear order is used directly, then many subtasks that could be
executed concurrently would be forced into a sequential schedule. To avoid
this conservatism, the accepting run is first decomposed into indexed subtask
occurrences. The complete order of these occurrences is then relaxed into a
partial-order representation that keeps only the essential temporal
constraints. Finally, an anytime branch-and-bound procedure is applied to
assign these subtasks to feasible agents or coalitions.

The section is organized as follows.
Subsection~\ref{subsec:po_subtask_extraction} introduces the extraction of
subtask occurrences from an accepting run.
Subsection~\ref{subsec:po_relaxation} defines the partial relations and
presents the computation of accepting posets.
Subsection~\ref{subsec:po_assignment} formulates the assignment problem over
a poset and presents an anytime search algorithm.
Subsection~\ref{subsec:po_analysis} summarizes the main correctness
properties.

\subsection{Automaton-Guided Subtask Extraction}
\label{subsec:po_subtask_extraction}

Consider the pruned B\"uchi automaton
$\mathcal{B}_{\varphi}^{-}=(Q^{-},Q^{-}_{0},\Sigma,\delta^{-},Q^{-}_{F})$
obtained from the global task formula $\varphi$. Since the task is assumed to
be co-safe, satisfaction can be certified by a finite prefix that reaches an
accepting state. Therefore, the decomposition step is based on accepting runs
of finite length. Let
$\rho=q_{0}q_{1}\cdots q_{L}$ be such a run, where
$q_{0}\in Q^{-}_{0}$ and $q_{L}\in Q^{-}_{F}$. Each transition along this run
encodes one logical progress step, and thus induces one subtask occurrence.

A careful distinction must be made between a subtask label and a subtask
occurrence. The same symbolic label may appear repeatedly in an accepting run,
but those appearances can have different logical meanings because they are
attached to different automaton states. For this reason, each occurrence is
indexed explicitly. In addition, besides the enabling label that triggers the
transition to the next state, one must also retain the self-loop requirement
that characterizes the propositions allowed or required while the automaton
remains at the current state. These two components are recorded separately.

\begin{definition}[Subtask occurrence]
\label{def:po_subtask_occurrence}
Let $\rho=q_{0}q_{1}\cdots q_{L}$ be an accepting run of
$\mathcal{B}_{\varphi}^{-}$. A subtask occurrence associated with the
$\ell$th transition is defined by the triple
\begin{equation}
\label{eq:po_subtask_occurrence}
\omega_{\ell} \triangleq (\ell,\sigma_{\ell},\sigma_{\ell}^{s}),
\end{equation}
where $\ell\in\{1,\cdots,L\}$ is the occurrence index,
$\sigma_{\ell}\subseteq\Sigma$ is a minimal enabling label for the transition
from $q_{\ell-1}$ to $q_{\ell}$, and
$\sigma_{\ell}^{s}\subseteq\Sigma$ is a minimal self-loop label at
$q_{\ell-1}$. More precisely,
\begin{equation}
\label{eq:po_minimal_transition}
q_{\ell}\in\delta^{-}(q_{\ell-1},\sigma_{\ell}), \qquad
q_{\ell}\notin\delta^{-}(q_{\ell-1},\bar{\sigma}),
\end{equation}
where $\bar{\sigma}\subsetneq\sigma_{\ell}$, and
\begin{equation}
\label{eq:po_minimal_selfloop}
q_{\ell-1}\in\delta^{-}(q_{\ell-1},\sigma_{\ell}^{s}), \qquad
q_{\ell-1}\notin\delta^{-}(q_{\ell-1},\bar{\sigma}^{s}),
\end{equation}
where $\bar{\sigma}^{s}\subsetneq\sigma_{\ell}^{s}$.
The set of all occurrences induced by $\rho$ is
\begin{equation}
\label{eq:po_occurrence_set}
\Omega_{\rho} \triangleq \{\omega_{1},\cdots,\omega_{L}\}.
\end{equation}
\end{definition}

The label $\sigma_{\ell}$ in \eqref{eq:po_minimal_transition} specifies the
smallest symbolic requirement that drives the automaton from
$q_{\ell-1}$ to $q_{\ell}$. The label $\sigma_{\ell}^{s}$ in
\eqref{eq:po_minimal_selfloop} identifies the smallest symbolic requirement
that keeps the automaton at $q_{\ell-1}$. The decomposition in
\eqref{eq:po_occurrence_set} therefore captures both progress and persistence.
This distinction is important later when temporal compatibility between
subtasks is analyzed.

\begin{remark}
\label{rem:po_occurrence_index}
The occurrence index in \eqref{eq:po_subtask_occurrence} is essential.
Two subtask occurrences with the same enabling label may still need to be
treated differently if they arise at different stages of the accepting run.
Discarding the index would collapse logically distinct requests into one
symbolic object and would generally destroy the temporal structure inherited
from the automaton.
\end{remark}

The set $\Omega_{\rho}$ obtained from \eqref{eq:po_occurrence_set} is ordered
implicitly by the run index. This order is always correct, but it can be
overly conservative for execution. In particular, it forbids any concurrency
between occurrences that appear in different positions of the run, even if the
automaton would still accept after reordering some of them. The next
subsection introduces a relaxed representation that preserves only the
necessary constraints.

\subsection{Relaxed Partial-Order Abstraction}
\label{subsec:po_relaxation}

A single binary precedence relation is not sufficient for the present
setting. The reason is that action propositions are not interpreted as
instantaneous events. Instead, a proposition remains true during the entire
execution interval of the associated action. As a result, two different types
of temporal information must be represented. The first indicates that one
subtask must start no later than another. The second indicates that a group of
subtasks cannot all be active at the same time. The original source material
captures this distinction through the relations
$\preceq_{\varphi}$ and $\neq_{\varphi}$, and the same notation is adopted
here for consistency.

\begin{definition}[Partial relations]
\label{def:po_partial_relations}
Given a decomposition $\Omega_{\rho}$, define the relation
$\preceq_{\varphi}\subseteq\Omega_{\rho}\times\Omega_{\rho}$ and the family
$\neq_{\varphi}\subseteq 2^{\Omega_{\rho}}$ as follows.

For $\omega_{h},\omega_{\ell}\in\Omega_{\rho}$,
\[
\omega_{h}\preceq_{\varphi}\omega_{\ell},
\]
means that $\omega_{h}$ must be started before $\omega_{\ell}$ is started.

For a subset $X\subseteq\Omega_{\rho}$ with $|X|\geq 2$,
\[
X\in\neq_{\varphi},
\]
means that the subtasks in $X$ cannot all be executed simultaneously.
\end{definition}

The relation $\preceq_{\varphi}$ describes ordering constraints among
subtasks. The family $\neq_{\varphi}$ describes concurrent exclusion
constraints. The two notions are fundamentally different. A precedence
constraint concerns the relative order of starting times, while an exclusion
constraint concerns the overlap of execution intervals. When actions have
nonzero duration, these two notions cannot be merged without losing relevant
temporal information.

\begin{definition}[Poset of subtasks]
\label{def:po_poset}
A poset associated with the accepting run $\rho$ is the triple
\begin{equation}
\label{eq:po_poset}
P_{\rho} \triangleq
(\Omega_{\rho},\preceq_{\varphi},\neq_{\varphi}).
\end{equation}
\end{definition}

The poset in \eqref{eq:po_poset} is more general than a classical partial
order, since the concurrent exclusion information is stored by a family of
subsets instead of a binary relation. This extension is necessary because a
forbidden simultaneous execution pattern may involve more than two subtask
occurrences. The representation remains compact, however, because it records
only those relations that are logically relevant to execution.

\begin{remark}
\label{rem:po_duration}
Most automata-based planning approaches for temporal logic treat the truth of
an action proposition as instantaneous
\cite{guo2015multi,tumova2016multi,kantaros2020stylus,
luo2021abstraction,luo2022temporal,jones2019scratchs,sahin2019multirobot}.
Under that convention, exact simultaneity is rarely meaningful, and the
relation $\neq_{\varphi}$ can often be ignored. In the present setting,
however, action propositions remain valid throughout their execution
intervals. Consequently, the distinction between
$\preceq_{\varphi}$ and $\neq_{\varphi}$ is necessary.
\end{remark}

The semantic meaning of a poset is given by the set of all timed words that
respect these relations. Let $t_{\ell}$ denote the starting time of
$\omega_{\ell}$, and let $d_{\ell}$ denote its duration. The execution of
$\omega_{\ell}$ then occupies the interval $[t_{\ell},\,t_{\ell}+d_{\ell})$.

\begin{definition}[Language of a poset]
\label{def:po_language}
Given a poset
$P_{\rho}=(\Omega_{\rho},\preceq_{\varphi},\neq_{\varphi})$,
its language is the set of finite timed words
\begin{equation}
\label{eq:po_timed_word}
W_{\rho} \triangleq
(t_{1},\omega_{1})(t_{2},\omega_{2})\cdots(t_{L},\omega_{L}),
\end{equation}
such that the following two conditions hold:
\begin{equation}
\label{eq:po_prec_constraint}
t_{h}\leq t_{\ell},
\end{equation}
whenever $\omega_{h}\preceq_{\varphi}\omega_{\ell}$, and
\begin{equation}
\label{eq:po_neq_constraint}
\forall X\in\neq_{\varphi}, \;
\exists \omega_{h},\omega_{\ell}\in X
\text{ such that }
t_{h}+d_{h}<t_{\ell},
\end{equation}
where $d_{h}$ and $d_{\ell}$ are the durations of
$\omega_{h}$ and $\omega_{\ell}$, respectively.
The language of $P_{\rho}$ is denoted by $L(P_{\rho})$.
\end{definition}

Equation~\eqref{eq:po_prec_constraint} encodes the precedence requirement.
Equation~\eqref{eq:po_neq_constraint} states that, for every subset in
$\neq_{\varphi}$, at least one member must finish before another member of
that same subset starts. Hence, the entire set cannot be active
simultaneously. This interpretation follows the original formulation and is
well suited for collaborative tasks with non-negligible duration.

A poset is useful only if it is still correct with respect to the original
task. This leads to the following notion.

\begin{definition}[Accepting poset]
\label{def:po_accepting}
A poset $P_{\rho}$ is called accepting if
\begin{equation}
\label{eq:po_accepting}
L(P_{\rho}) \subseteq L_{\varphi},
\end{equation}
where $L_{\varphi}$ is the language of the global task formula.
\end{definition}

The computation starts from the fully ordered version of the accepting run and
then removes redundant relations. More specifically, let
\begin{equation}
\label{eq:po_full_order}
\mathbb{F} \triangleq
\{(\omega_{h},\omega_{\ell})\in\Omega_{\rho}\times\Omega_{\rho}
\mid h<\ell\}.
\end{equation}
The initial poset is taken as
$(\Omega_{\rho},\mathbb{F},2^{\Omega_{\rho}})$.
This choice is maximally conservative. It is then relaxed by checking
acceptance after local perturbations of the timed word associated with the
accepting run. If a swap of two adjacent occurrences still yields an accepted
word, the corresponding order relation is not essential. If a simultaneous
realization of a subset still yields an accepted word, that subset need not
remain in $\neq_{\varphi}$.

\begin{algorithm}[t!]
\caption{$\texttt{ComputePoset}(\mathcal{B}_{\varphi}^{-},t_{0})$}
\label{alg:po_extract}
\SetKwInOut{Input}{Input}
\SetKwInOut{Output}{Output}
\Input{Pruned automaton $\mathcal{B}_{\varphi}^{-}$ and time budget $t_{0}$.}
\Output{A set of accepting posets $\mathcal{P}_{\varphi}$ and their language
cover $\mathcal{L}_{\varphi}$.}

Choose $q_{0}\in Q_{0}^{-}$ and $q_{f}\in Q_{F}^{-}$\;
Initialize $\mathcal{P}_{\varphi}\leftarrow\emptyset$ and
$\mathcal{L}_{\varphi}\leftarrow\emptyset$\;
Start a modified DFS on $\mathcal{B}_{\varphi}^{-}$ to generate an accepting
run $\rho=q_{0}q_{1}\cdots q_{L}$\;

\While{$\mathrm{time}<t_{0}$}{
  Compute $\Omega_{\rho}$ from \eqref{eq:po_occurrence_set}\;
  Construct the canonical word
  $W_{\rho}\leftarrow(\omega_{1}\omega_{2}\cdots\omega_{L})$\;

  \If{$W_{\rho}\notin\mathcal{L}_{\varphi}$}{
    Set $\preceq_{\varphi}\leftarrow\mathbb{F}$ by
    \eqref{eq:po_full_order}\;
    Set $\neq_{\varphi}\leftarrow 2^{\Omega_{\rho}}$\;
    Set $P_{\rho}\leftarrow(\Omega_{\rho},\preceq_{\varphi},\neq_{\varphi})$\;
    Set $L(P_{\rho})\leftarrow\{W_{\rho}\}$ and
    $\mathrm{Que}\leftarrow\{W_{\rho}\}$\;
    Set $I_{\mathrm{acc}}\leftarrow\emptyset$ and
    $I_{\mathrm{rej}}\leftarrow\emptyset$\;

    \While{$\mathrm{Que}\neq\emptyset$}{
      Pop one word $W=(\eta_{1}\eta_{2}\cdots\eta_{L})$ from $\mathrm{Que}$\;

      \For{$i=1$ \KwTo $L-1$}{
        Set $\widetilde{W}\leftarrow
        (\eta_{1}\cdots\eta_{i+1}\eta_{i}\cdots\eta_{L})$\;

        \If{$\widetilde{W}$ is accepted by $\mathcal{B}_{\varphi}^{-}$}{
          Add $(\eta_{i},\eta_{i+1})$ to $I_{\mathrm{acc}}$\;
          Add $\widetilde{W}$ to $\mathrm{Que}$ if unseen\;
          Add $\widetilde{W}$ to $L(P_{\rho})$ if unseen\;
        }
        \Else{
          Add $(\eta_{i},\eta_{i+1})$ to $I_{\mathrm{rej}}$\;
        }
      }

      Update
      $\preceq_{\varphi}\leftarrow
      \preceq_{\varphi}\setminus
      (I_{\mathrm{acc}}\setminus I_{\mathrm{rej}})$\;

      \ForEach{$X\subseteq\Omega_{\rho}$ such that $X\in\neq_{\varphi}$}{
        Construct $\widehat{W}$ by replacing all occurrences in $X$
        within $W$ by their simultaneous union label\;

        \If{$\widehat{W}$ is accepted by $\mathcal{B}_{\varphi}^{-}$}{
          Update
          $\neq_{\varphi}\leftarrow\neq_{\varphi}\setminus\{X\}$\;
        }
      }
    }

    Update $P_{\rho}\leftarrow
    (\Omega_{\rho},\preceq_{\varphi},\neq_{\varphi})$\;
    Add $P_{\rho}$ to $\mathcal{P}_{\varphi}$\;
    Update
    $\mathcal{L}_{\varphi}\leftarrow
    \mathcal{L}_{\varphi}\cup L(P_{\rho})$\;
  }

  Continue the modified DFS until the next accepting run is found\;
}

\Return{$\mathcal{P}_{\varphi},\mathcal{L}_{\varphi}$}\;
\end{algorithm}

Algorithm~\ref{alg:po_extract} makes the relaxation process explicit.
The initial relation $\preceq_{\varphi}=\mathbb{F}$ imposes the exact order
of the accepting run, and $\neq_{\varphi}=2^{\Omega_{\rho}}$ prohibits any
simultaneous execution. The algorithm then reduces these relations by direct
acceptance tests on $\mathcal{B}_{\varphi}^{-}$. A precedence edge is removed
only if the corresponding adjacent swap is accepted. A subset is removed from
$\neq_{\varphi}$ only if its simultaneous realization is accepted. Therefore,
every relaxation is certified rather than guessed.

\begin{theorem}
\label{thm:po_accepting_poset}
Any poset returned by Algorithm~\ref{alg:po_extract} is accepting.
\end{theorem}

\begin{proof}
Consider a poset $P_{\rho}$ returned by
Algorithm~\ref{alg:po_extract}. The initial word associated with the
accepting run $\rho$ is accepted by $\mathcal{B}_{\varphi}^{-}$ by
construction. During the algorithm, a precedence relation is removed only
after the swapped word has been checked and accepted by
$\mathcal{B}_{\varphi}^{-}$. Likewise, a subset is removed from
$\neq_{\varphi}$ only after the corresponding simultaneous realization has
been checked and accepted. Hence, every timed word added to $L(P_{\rho})$
corresponds to accepted automaton behavior. It follows that
$L(P_{\rho})\subseteq L_{\varphi}$, and thus $P_{\rho}$ is accepting.
\end{proof}

\begin{remark}
\label{rem:po_compare}
The abstraction above includes decompositions based on decomposable states as
a special case \cite{schillinger2018simultaneous}. The essential difference
is that the present construction may reveal concurrency not only between
different segments of an accepting run, but also within a segment, provided
that the corresponding acceptance tests succeed. This additional flexibility
is important for minimizing the completion time of collaborative tasks.
\end{remark}

\subsection{Partial-Order-Constrained Task Assignment}
\label{subsec:po_assignment}

Once an accepting poset has been computed, the global temporal task is
reduced to an assignment problem over subtask occurrences. The logical
constraints are already encoded by $P_{\rho}$. What remains is to decide
which agents or coalitions should execute each occurrence, and in what order,
so that the resulting timed execution respects the poset and minimizes the
overall completion time.

Let $\mathcal{N}\triangleq\{1,\cdots,N\}$ be the agent index set.
For each occurrence $\omega\in\Omega_{\rho}$, define
$\Gamma(\omega)\subseteq 2^{\mathcal{N}}$ as the collection of feasible
executing coalitions. If $I\in\Gamma(\omega)$ and $|I|=1$, then $\omega$
corresponds to a local task. If $|I|>1$, then $\omega$ requires direct
collaboration among multiple agents. This notation is consistent with the
coalition-based task model introduced earlier in the chapter.

\begin{definition}[Poset graph]
\label{def:po_graph}
Given the poset
$P_{\rho}=(\Omega_{\rho},\preceq_{\varphi},\neq_{\varphi})$,
its associated poset graph is the directed graph
\begin{equation}
\label{eq:po_graph}
G_{\rho} \triangleq (\Omega_{\rho},E_{\rho}),
\end{equation}
where $(\omega_{h},\omega_{\ell})\in E_{\rho}$ if
$\omega_{h}\preceq_{\varphi}\omega_{\ell}$ and there exists no
$\omega_{m}\in\Omega_{\rho}$ such that
$\omega_{h}\preceq_{\varphi}\omega_{m}\preceq_{\varphi}\omega_{\ell}$ with
$\omega_{m}\neq\omega_{h}$ and $\omega_{m}\neq\omega_{\ell}$.
For any $\omega\in\Omega_{\rho}$, the predecessor set of $\omega$ in
$G_{\rho}$ is denoted by $\mathrm{Pred}(\omega)$.
\end{definition}

The graph in \eqref{eq:po_graph} removes all transitive edges and therefore
provides the minimal causal structure required for node expansion in the
assignment algorithm. This avoids unnecessary constraints during search while
preserving all logical dependencies encoded by $\preceq_{\varphi}$.

\begin{problem}[Poset-constrained task assignment]
\label{prob:po_assignment}
Given an accepting poset
$P_{\rho}=(\Omega_{\rho},\preceq_{\varphi},\neq_{\varphi})$,
determine for each $\omega\in\Omega_{\rho}$ a feasible coalition
$I_{\omega}\in\Gamma(\omega)$ and a start time $t_{\omega}$ such that all
constraints in \eqref{eq:po_prec_constraint} and
\eqref{eq:po_neq_constraint} are satisfied, and the makespan is minimized.
\end{problem}

Let $\mathfrak{J}(P_{\rho})$ denote the set of feasible assignments for
Problem~\ref{prob:po_assignment}. The optimal assignment is
\begin{equation}
\label{eq:po_opt_assignment}
J^{\star} \triangleq
\arg\min_{J\in\mathfrak{J}(P_{\rho})} T(J),
\end{equation}
where $T(J)$ is the completion time of assignment $J$.

Problem~\ref{prob:po_assignment} is combinatorial even without the temporal
logic structure, since it contains routing, precedence-constrained
scheduling, and coalition formation as special cases
\cite{hochba1997approximation,morrison2016branch}. For this reason, an
anytime branch-and-bound procedure is adopted. The search tree is built over
partial assignments, and only those subtasks whose predecessors have already
been assigned are eligible for expansion.

\begin{definition}[Search node]
\label{def:po_node}
A search node is a tuple
\begin{equation}
\label{eq:po_search_node}
\nu \triangleq (\tau_{1},\cdots,\tau_{N}),
\end{equation}
where $\tau_{n}$ is the ordered sequence of occurrences assigned to agent
$n\in\mathcal{N}$.
The set of assigned subtasks at node $\nu$ is
\begin{equation}
\label{eq:po_assigned_set}
\Omega_{\nu} \triangleq \bigcup_{n\in\mathcal{N}} \tau_{n}.
\end{equation}
The set of currently expandable occurrences is
\begin{equation}
\label{eq:po_expandable}
\mathrm{Next}(\nu) \triangleq
\left\{
\omega\in\Omega_{\rho}\setminus\Omega_{\nu} \;\middle|\;
\mathrm{Pred}(\omega)\subseteq\Omega_{\nu}
\right\}.
\end{equation}
\end{definition}

Equation~\eqref{eq:po_expandable} is the key admissibility condition for node
expansion. It ensures that no occurrence is assigned before all of its causal
predecessors have already appeared in the partial assignment. Hence, logical
correctness is preserved directly at the search-tree level.

The branch-and-bound method maintains one incumbent feasible assignment and
its current makespan. For every node, an upper bound is computed by greedily
completing the current partial assignment. A lower bound is computed from a
relaxation of the remaining problem. If the lower bound of a child node is no
smaller than the incumbent makespan, that branch cannot improve the current
solution and is discarded.

\begin{algorithm}[t!]
\caption{$\texttt{BnB\_Assign}(P_{\rho},t_{0})$}
\label{alg:po_bnb}
\SetKwInOut{Input}{Input}
\SetKwInOut{Output}{Output}
\Input{Accepting poset $P_{\rho}$ and time budget $t_{0}$.}
\Output{Best assignment $J^{\star}$ and best makespan $T^{\star}$.}

Initialize the root node
$\nu_{0}\leftarrow(\emptyset,\cdots,\emptyset)$\;
Initialize the active queue
$\mathcal{Q}\leftarrow\{(\nu_{0},0)\}$\;
Set $J^{\star}\leftarrow\emptyset$ and $T^{\star}\leftarrow+\infty$\;

\While{$\mathcal{Q}\neq\emptyset$ and $\mathrm{time}<t_{0}$}{
  Extract from $\mathcal{Q}$ one node $\nu$ with minimum lower bound\;

  Compute a greedy feasible completion
  $(\widehat{J}_{\nu},\widehat{T}_{\nu})$ of $\nu$\;

  \If{$\widehat{T}_{\nu}<T^{\star}$}{
    Set $J^{\star}\leftarrow\widehat{J}_{\nu}$\;
    Set $T^{\star}\leftarrow\widehat{T}_{\nu}$\;
  }

  Compute $\mathrm{Next}(\nu)$ from \eqref{eq:po_expandable}\;

  \ForEach{$\omega\in\mathrm{Next}(\nu)$}{
    \ForEach{$I\in\Gamma(\omega)$}{
      Construct the child node
      $\nu^{+}=\mathrm{Append}(\nu,\omega,I)$\;
      Compute the lower bound $\underline{T}(\nu^{+})$\;

      \If{$\underline{T}(\nu^{+})<T^{\star}$}{
        Insert $(\nu^{+},\underline{T}(\nu^{+}))$ into $\mathcal{Q}$\;
      }
    }
  }
}

\Return{$J^{\star},T^{\star}$}\;
\end{algorithm}

Algorithm~\ref{alg:po_bnb} is anytime in the usual sense.
Once the first feasible completion is found, an incumbent assignment exists.
As the search continues, the incumbent value can only improve. Therefore, the
algorithm can be terminated at any time and still returns the best feasible
solution computed so far. This property is particularly useful in large-scale
multi-robot missions, where a near-optimal solution obtained quickly may be
more valuable than an exact optimum obtained after a long delay.

\begin{theorem}
\label{thm:po_assignment_sound}
Any complete assignment returned by Algorithm~\ref{alg:po_bnb} satisfies the
poset constraints of $P_{\rho}$.
\end{theorem}

\begin{proof}
A node is expanded only with occurrences in $\mathrm{Next}(\nu)$, as defined
in \eqref{eq:po_expandable}. Therefore, every assigned occurrence appears
after all of its predecessors in the poset graph have already been assigned.
This guarantees that the precedence constraints encoded by
$\preceq_{\varphi}$ are respected. In addition, the timing synthesis step used
to evaluate a complete assignment accepts only schedules satisfying
\eqref{eq:po_neq_constraint}. Hence, every complete assignment returned by
Algorithm~\ref{alg:po_bnb} satisfies both
$\preceq_{\varphi}$ and $\neq_{\varphi}$.
\end{proof}

\subsection{Correctness and Discussion}
\label{subsec:po_analysis}

The two algorithms developed above form a coherent top-down synthesis layer.
Algorithm~\ref{alg:po_extract} computes accepting posets directly from the
task automaton. Algorithm~\ref{alg:po_bnb} searches for a low-makespan
assignment that respects one such poset. Since the first stage guarantees
logical admissibility and the second stage preserves the poset constraints,
their composition yields a sound solution method for the global temporal task.

\begin{theorem}[Soundness of the synthesis framework]
\label{thm:po_soundness}
Let $P_{\rho}$ be an accepting poset returned by
Algorithm~\ref{alg:po_extract}. Let $J^{\star}$ be a complete assignment
returned by Algorithm~\ref{alg:po_bnb} for $P_{\rho}$. Then the timed
execution induced by $J^{\star}$ satisfies the global task formula
$\varphi$.
\end{theorem}

\begin{proof}
By Theorem~\ref{thm:po_accepting_poset}, the poset $P_{\rho}$ satisfies
$L(P_{\rho})\subseteq L_{\varphi}$. By
Theorem~\ref{thm:po_assignment_sound}, the timed execution induced by
$J^{\star}$ belongs to $L(P_{\rho})$. Hence it also belongs to
$L_{\varphi}$. Therefore, the execution generated by $J^{\star}$ satisfies
the original global task formula $\varphi$.
\end{proof}

\begin{remark}
\label{rem:po_complexity}
The worst-case complexity of the overall procedure remains exponential in the
number of subtask occurrences and in the number of agents. This is expected,
since the assignment problem already contains NP-hard subproblems. The main
algorithmic advantage is not polynomial complexity, but rather the ability to
extract useful concurrency from the automaton and to maintain an incumbent
solution throughout the search. This combination is particularly attractive
for time-critical robotic applications.
\end{remark}

The present section establishes the logical and combinatorial basis of the
top-down synthesis framework. The partial-order abstraction exposes the
execution flexibility hidden in an accepting automaton run, while the
branch-and-bound assignment algorithm converts this flexibility into a concrete
multi-robot schedule. The next section can build on this representation to
address subsequent synthesis components without returning to the full global
automaton.

%==============================
%==============================
\section{Hierarchical Coordination of Large-Scale Fleets under Temporal Tasks}
\label{sec:hierarchical_coordination}

Large-scale temporal-task coordination for heterogeneous fleets faces two
coupled sources of complexity. The first source is logical. Temporal missions
must be decomposed into admissible task sequences that satisfy ordering,
concurrency, and completion requirements. The second source is combinatorial.
Concrete robots with different capabilities must be grouped into teams and
dispatched over the workspace with bounded response time. Direct formulations
that synchronize temporal automata with all robot models can provide strong
correctness guarantees, but their size grows rapidly with the fleet cardinality
and the mission complexity
\cite{schillinger2018simultaneous,kantaros2020stylus,luo2021temporal,
leahy2021scalable}. Recent advances have improved scalability by sampling,
decision-tree expansion, partial-order reduction, and hierarchical task
abstractions
\cite{chen2025real,liu2024time,luo2025simultaneous,gosrich2025online,
luo2025hulk}. Nevertheless, large fleets still require a sharper separation
between temporal decomposition, team formation, and execution-level
realization.

This section develops such a separation. The emphasis is the coordination
backbone that connects temporal specifications to executable team plans. The
presentation deliberately suppresses interface-level supervision and keeps
re-optimization details to a minimum. Those aspects are deferred to the next
section. The resulting architecture contains three tightly coupled layers. The
fleet level reasons over temporal tasks and constructs abstract team plans. The
intermediate level instantiates these abstract plans by assigning concrete
robots to capability profiles. The execution level converts each team plan into
robot-level plans $\tau_i$, which have been defined previously as timed action
sequences.

\subsection{Hierarchical coordination model}
\label{subsec:hierarchical_model}

Let $\mathcal{N}\triangleq\{1,\cdots,N\}$ denote the fleet and let
$\Phi_t\triangleq\{\varphi_1,\cdots,\varphi_M\}$ denote the set of temporal
missions known at time $t$. Each mission $\varphi_\ell$ is interpreted as a
syntactically co-safe temporal formula over collaborative tasks. This modeling
choice is natural for finite-horizon missions in which completion can be
certified after a finite trace, and it is consistent with automata-based
planning for temporal logic \cite{baier2008principles}. The task abstraction
must capture both capability requirements and the spatial support over which
collaboration is executed.

\begin{definition}[Collaborative task]
\label{def:collaborative_task}
A collaborative task is a tuple
\[
\omega \triangleq \big(S_\omega,\eta_\omega,\mathcal{J}_\omega\big),
\]
where $S_\omega\subseteq\mathcal{W}$ is the execution region,
$\eta_\omega$ is a duration map, and
$\mathcal{J}_\omega\triangleq\{(n_j,a_j,s_j)\}_{j=1}^{J_\omega}$ is the set
of subtasks. For each subtask, $n_j$ is the required number of robots,
$a_j$ is the required action, and $s_j\in S_\omega$ is the service location.
\end{definition}

Definition~\ref{def:collaborative_task} extends purely counting-based task
descriptions by retaining the spatial support and the dependence of execution
time on the assigned team. This distinction is important in large-scale fleets,
because the same cardinality requirement may induce very different response
times when the available robots differ in location and capability
\cite{sahin2019multirobot,cardona2024planning}. It is also convenient for
hierarchical coordination, since the fleet level only needs abstract task
information, while the execution level resolves robot trajectories.

The intermediate objects of the hierarchy are task sequences assigned to teams
and the associated capacity profiles.

\begin{definition}[Abstract team plan and capacity profile]
\label{def:team_profile}
For team index $k\in\mathcal{K}$, an abstract team plan is a finite sequence
\[
\Gamma_k \triangleq \omega_k^1\omega_k^2\cdots\omega_k^{L_k}.
\]
The capacity profile induced by $\Gamma_k$ is the map
$\beta_k:\mathcal{A}\rightarrow\mathbb{N}$ defined by
\begin{equation}
\label{eq:cap_profile_ch4}
\beta_k(a) \triangleq
\max_{\omega\in\Gamma_k}
\max_{\substack{(n_j,a_j,s_j)\in\mathcal{J}_\omega\\ a_j=a}}
n_j,
\end{equation}
where $\beta_k(a)$ is taken as zero when action $a$ does not appear in
$\Gamma_k$.
\end{definition}

Equation~\eqref{eq:cap_profile_ch4} records the largest instantaneous demand
for each action along the team sequence. It deliberately does not identify
which robots realize that demand. This abstraction is the key mechanism that
decouples fleet-level temporal decomposition from concrete robot assignment.

\subsection{Automaton-guided fleet-level assignment}
\label{subsec:automaton_assignment}

For each mission $\varphi_\ell\in\Phi_t$, let
\[
\mathcal{B}_{\varphi_\ell}
\triangleq
\big(Q_\ell,Q_\ell^0,\Sigma_\ell,\delta_\ell,F_\ell\big)
\]
denote the corresponding nondeterministic B\"uchi or finite-word acceptor,
depending on the exact temporal fragment introduced previously. The fleet-level
coordination problem is to construct a set of abstract team plans
$\{\Gamma_k\}_{k\in\mathcal{K}}$ and an admissible number of teams $K$ such
that the induced trace satisfies all mission automata and the aggregate
capability demand remains feasible.

\begin{problem}
\label{prob:fleet_level_coordination}
Given the mission set $\Phi_t$ and the fleet capability model
$\{\mathcal{A}_i\}_{i\in\mathcal{N}}$, determine the number of teams $K$, the
abstract team plans $\{\Gamma_k\}_{k\in\mathcal{K}}$, and the associated
capacity profiles $\{\beta_k\}_{k\in\mathcal{K}}$ such that all missions are
accepted by their automata and the fleet response time is minimized.
\end{problem}

A direct synchronization of all robot models with all mission automata is
avoided. Instead, the search is carried out over partial team plans together
with the currently reachable automaton states.

\begin{definition}[Assignment node]
\label{def:assignment_node}
Let $\mathcal{M}\triangleq\{1,\cdots,M\}$. An assignment node is a tuple
\[
\nu \triangleq
\big(\{\Gamma_k\}_{k\in\mathcal{K}},\{\widehat{Q}_\ell\}_{\ell\in\mathcal{M}}\big),
\]
where $\Gamma_k$ is the current task sequence of team $k$ and
$\widehat{Q}_\ell\subseteq Q_\ell$ is the set of automaton states reachable in
$\mathcal{B}_{\varphi_\ell}$ after executing the task symbols already assigned
in the node.
\end{definition}

The root node contains empty team plans and the initial reachable-state sets.
Each expansion appends one enabled task symbol to one existing team, or to a
newly created team, and updates the reachable-state sets accordingly. A task
symbol is enabled at node $\nu$ if it appears on at least one outgoing
transition of some state in $\widehat{Q}_\ell$ for at least one mission
$\ell\in\mathcal{M}$. This rule allows temporal decomposition and team
assignment to proceed simultaneously.

To rank candidate nodes, a multi-term value function is introduced:
\begin{equation}
\label{eq:node_value_ch4}
\chi(\nu) \triangleq
\max_{k\in\mathcal{K}} T_k
+ \lambda_c \sum_{k\in\mathcal{K}} C_k
+ \lambda_p \sum_{\ell\in\mathcal{M}}
w_\ell \min_{q\in\widehat{Q}_\ell}\psi_\ell(q,F_\ell),
\end{equation}
where $T_k$ is the estimated completion time of $\Gamma_k$, $C_k$ is the
estimated travel-and-service cost of $\Gamma_k$, $w_\ell$ is the prescribed
importance weight of mission $\varphi_\ell$, and
$\psi_\ell(q,F_\ell)$ is the shortest graph distance from $q$ to the accepting
set $F_\ell$,
where $\lambda_c>0$ and $\lambda_p>0$ are regularization weights.

Equation~\eqref{eq:node_value_ch4} has a clear interpretation. The first term
controls the makespan of the current partial assignment. The second term
penalizes unnecessarily expensive team sequences. The third term favors nodes
that are logically closer to acceptance. A purely makespan-driven expansion is
therefore avoided, which improves search stability when many nodes share the
same temporal depth.

Every node must also satisfy a fleet-level capability test. Using the capacity
profiles defined in Definition~\ref{def:team_profile}, feasibility requires
\begin{equation}
\label{eq:global_capacity_ch4}
\sum_{k\in\mathcal{K}} \beta_k(a)
\leq
\sum_{i\in\mathcal{N}}\mathds{1}\big(a\in\mathcal{A}_i\big),
\quad \forall a\in\mathcal{A},
\end{equation}
where the left-hand side is the total demand for action $a$ induced by the
current abstract team plans, and the right-hand side is the number of robots in
the fleet that can perform action $a$.

A further reduction is obtained through dominance pruning. For each node,
define the profile
\begin{equation}
\label{eq:node_profile_ch4}
\rho(\nu) \triangleq
\Big[
\{T_k\}_{k\in\mathcal{K}},
\{C_k\}_{k\in\mathcal{K}},
\Big\{
\min_{q\in\widehat{Q}_\ell}\psi_\ell(q,F_\ell)
\Big\}_{\ell\in\mathcal{M}}
\Big].
\end{equation}
Node $\nu_1$ dominates node $\nu_2$ when
$\rho(\nu_1)\leq\rho(\nu_2)$ componentwise and at least one inequality is
strict. Only non-dominated nodes are retained on the frontier. The search thus
progresses over a compressed representation of the assignment space instead of a
full synchronous product.

\begin{algorithm}[t]
\caption{Automaton-guided Hierarchical Task Assignment}
\label{alg:hierarchical_assignment}
\SetAlgoLined
\KwIn{Mission set $\Phi_t$, automata $\{\mathcal{B}_{\varphi_\ell}\}$,
fleet $\mathcal{N}$, horizon $H$.}
\KwOut{Number of teams $K$, abstract team plans $\{\Gamma_k\}$,
capacity profiles $\{\beta_k\}$.}

Initialize the root node
$\nu_0 \gets (\{\emptyset\},\{Q_\ell^0\}_{\ell\in\mathcal{M}})$\;
Initialize the frontier $\mathscr{F}\gets\{\nu_0\}$\;

\While{the budget is not exhausted}{
  Select a batch of frontier nodes with smallest
  $\chi(\nu)$ in~\eqref{eq:node_value_ch4}\;

  \ForEach{selected node $\nu$}{
    compute the enabled task-symbol set from
    $\{\widehat{Q}_\ell\}_{\ell\in\mathcal{M}}$\;

    \ForEach{enabled symbol $\omega$}{
      \ForEach{existing team $k$ and the option of creating one new team}{
        append $\omega$ to $\Gamma_k$\;
        update the reachable-state sets
        $\{\widehat{Q}_\ell\}_{\ell\in\mathcal{M}}$\;
        update $T_k$, $C_k$, and $\beta_k(\cdot)$\;
        \lIf{\eqref{eq:global_capacity_ch4} is violated}{discard the node}
        \lIf{the number of newly committed tasks exceeds $H$}{pause expansion}
      }
    }
  }

  prune dominated nodes using~\eqref{eq:node_profile_ch4}\;
  update the frontier with the non-dominated surviving nodes\;
}

return the complete frontier node with minimum
$\chi(\nu)$ in~\eqref{eq:node_value_ch4}\;
\end{algorithm}

Algorithm~\ref{alg:hierarchical_assignment} summarizes the fleet-level
backbone. The role of the horizon $H$ is purely computational in the present
section. It bounds the number of future task commitments in one optimization
cycle and prevents the assignment layer from becoming excessively myopic or
excessively rigid. Detailed event-triggered replanning policies are omitted
here and are addressed separately later.

The logical correctness of the fleet-level layer follows directly from the node
semantics and the acceptance condition of the mission automata.

\begin{proposition}
\label{prop:fleet_correctness}
Assume that the search horizon contains all admissible task symbols for the
current planning instance. If Algorithm~\ref{alg:hierarchical_assignment}
returns a complete feasible node, then the resulting abstract team plans
$\{\Gamma_k\}_{k\in\mathcal{K}}$ satisfy all mission formulas and all aggregate
capability constraints in~\eqref{eq:global_capacity_ch4}.
\end{proposition}

The proof is immediate. Each node expansion follows valid automaton
transitions, and completeness requires that every mission automaton reaches an
accepting set. Capability feasibility is enforced explicitly through
\eqref{eq:global_capacity_ch4}. Hence any complete feasible node corresponds to
a temporally valid fleet-level decomposition.

\subsection{Capacity-aware team formation}
\label{subsec:capacity_team_formation}

The fleet-level layer determines what each team must be able to do, but it does
not determine which robots belong to which team. The next layer instantiates
each abstract team plan by assigning concrete robots to the induced capacity
profile. This assignment must preserve disjoint team membership and may include
limited redundancy to absorb uncertainty in workload or travel time.

\begin{problem}
\label{prob:team_formation}
Given the abstract team plans $\{\Gamma_k\}_{k\in\mathcal{K}}$ returned by
Algorithm~\ref{alg:hierarchical_assignment}, determine the concrete teams
$\{\mathcal{N}_k\}_{k\in\mathcal{K}}$ such that each team can realize the
capacity profile induced by its plan and the worst team completion time is
minimized.
\end{problem}

Introduce binary variables $b_{ik}\in\{0,1\}$, where $b_{ik}=1$ means that
robot $i$ is assigned to team $k$. Let $\alpha_a\geq 1$ denote the redundancy
margin associated with action $a$. The formation constraints are
\begin{equation}
\label{eq:team_formation_constraints_ch4}
\beta_k(a)
\leq
\sum_{i\in\mathcal{N}} b_{ik}\mathds{1}\big(a\in\mathcal{A}_i\big)
\leq
\alpha_a\,\beta_k(a),
\quad
\sum_{k\in\mathcal{K}} b_{ik}\leq 1,
\end{equation}
where the first inequality enforces the minimum and maximum staffing level of
team $k$ for action $a$, and the second inequality guarantees disjoint team
membership.

The lower bound in~\eqref{eq:team_formation_constraints_ch4} guarantees that a
team can execute every task in its abstract plan. The upper bound prevents
excessive concentration of resources in one team. This redundancy window is
especially useful when some tasks contain partially revealed subtasks or when
travel-time estimates are conservative. The abstraction remains much smaller
than direct robot-to-subtask assignment, because the binary variables appear at
the level of teams rather than at the level of all subtasks.

Let $S_k^1$ denote the first region in $\Gamma_k$. For each robot $i$ and team
$k$, define the estimated arrival time
$t_{ik}^{\mathrm{arr}}\triangleq
\widehat{t}_i+T_{\mathrm{nav}}(\widehat{x}_i,S_k^1)$, where $\widehat{t}_i$
and $\widehat{x}_i$ are the current availability time and current position of
robot $i$. The completion cost of team $k$ is then modeled by
\begin{equation}
\label{eq:team_completion_cost_ch4}
J_k \triangleq
\max_{i\,:\,b_{ik}=1} t_{ik}^{\mathrm{arr}}
+
\sum_{\omega\in\Gamma_k} T_{\mathrm{exec}}(\omega),
\end{equation}
where the first term captures the synchronization delay of the slowest assigned
robot and the second term is the total service time of the team plan.

The team-formation layer minimizes the worst team completion time, namely
$\min \max_{k\in\mathcal{K}} J_k$, subject to
\eqref{eq:team_formation_constraints_ch4}. Once the binary solution is
obtained, the concrete team is recovered by
$\mathcal{N}_k\triangleq\{i\in\mathcal{N}\mid b_{ik}=1\}$. The outcome of this
layer is therefore a physically realizable partition of the fleet that remains
fully consistent with the abstract temporal decomposition.

\subsection{Execution-level realization of team plans}
\label{subsec:execution_realization}

After concrete teams have been formed, each abstract team plan is converted into
a rollout over regions and tasks:
\begin{equation}
\label{eq:team_rollout_ch4}
\Xi_k \triangleq
(S_k^1,\omega_k^1)\cdots(S_k^{L_k},\omega_k^{L_k}),
\end{equation}
where $S_k^\ell$ is the execution region of task $\omega_k^\ell$.
Equation~\eqref{eq:team_rollout_ch4} is the interface between the upper
coordination layers and the robot-level layer. It specifies what each team must
execute and in which temporal order, but it does not prescribe detailed robot
motions.

The execution layer resolves this remaining degree of freedom inside each team.
Given $\Xi_k$ and the team $\mathcal{N}_k$, a local coordinator generates the
robot-level plans $\tau_i$ and the associated trajectories. At this level,
motion feasibility, synchronization among robots of the same team, and the
internal assignment of subtasks are handled locally. The exact local mechanism
depends on the task family. Static known subtasks lead to routing and
trajectory-optimization problems. Partially revealed subtasks require
exploration-aware insertion policies. Dynamic subtasks require coalition updates
inside the team. These mechanisms are essential for realization, yet they do not
alter the logic of the hierarchy itself. Their role is to refine the team plan,
not to redefine it.

This separation yields the main structural advantage of the proposed framework.
Temporal satisfaction is decided at the fleet level through automaton-guided
search. Concrete resource instantiation is handled at the team-formation level
through a compact capacity-based assignment. Motion planning and local
collaboration are confined to the execution layer. Compared with monolithic
temporal-task formulations, the computational burden is therefore distributed
over smaller and better structured subproblems
\cite{kantaros2020stylus,leahy2021scalable,chen2025real,
luo2025simultaneous,gosrich2025online,luo2025hulk}. This hierarchical
decomposition is the essential reason why large fleets can be coordinated under
rich temporal tasks without explicitly constructing the full synchronous product
of all robot models and all mission automata.

%==============================
%==============================
\section{Online Adaptation for Continual Tasks}
\label{sec:online_continual_tasks}

The previous sections considered collaborative task planning for a given set of
high-level specifications. In continual operation, this assumption no longer
holds. New tasks may be released while the robots are executing previously
assigned subtasks, and the resulting plan must be updated without discarding
valid progress that has already been achieved. Direct recomputation from the
conjunction of all active formulas is generally undesirable, since the
associated automaton construction and the subsequent multi-robot search may
become prohibitively expensive for long temporal formulas
\cite{baier2008principles,schillinger2018simultaneous,
kantaros2020stylus,luo2021temporal}. For continual tasks, a more suitable
strategy is to update the high-level partial-order structure online and then
re-solve only a finite planning window over the updated structure
\cite{liu2024fast,liu2024time}.

This section develops such an online adaptation mechanism. The first component
is an incremental product of relaxed partially ordered sets, abbreviated as
R-posets, by which a newly released temporal task is merged into the current
global partial order. The second component is a receding-horizon planning
module, inherited from the previous section, that is repeatedly solved on a
finite subset of currently relevant subtasks. The resulting architecture
preserves the correctness of the high-level task logic while maintaining
computational tractability during long-duration execution.

\subsection{Continual task model and planning state}
\label{subsec:continual_model}

Let $t \geq 0$ denote the current planning time. The set of active continual
task formulas is denoted by
$\Phi_t \triangleq \{\varphi_t^1,\cdots,\varphi_t^{m_t}\}$, where each
$\varphi_t^\ell$ is an sc-LTL formula over the atomic propositions introduced
earlier in the chapter. The currently maintained global R-poset is denoted by
\[
\mathcal{P}_t \triangleq (\Omega_t,\preceq_t,\neq_t),
\]
where $\Omega_t$ is the set of active subtasks, $\preceq_t$ is the partial
precedence relation, and $\neq_t$ is the conflict relation. As in the previous
sections, a subtask in $\Omega_t$ may represent either an individual action or
a collaborative action that requires multiple robots.

The online update should account for the fact that part of the current poset
may already have been executed. To this end, let
$\Omega_t^{\mathrm{fin}} \subseteq \Omega_t$ be the set of finished subtasks,
let $\Omega_t^{\mathrm{run}} \subseteq \Omega_t \setminus
\Omega_t^{\mathrm{fin}}$ be the set of subtasks currently under execution, and
let $\Omega_t^{\mathrm{rem}} \triangleq \Omega_t \setminus
\Omega_t^{\mathrm{fin}}$ be the set of remaining subtasks. Moreover, let
$\mathbf{w}_t^{\mathrm{exe}}$ denote the finite execution word that has already
been generated by the system up to time $t$.

\begin{definition}[Continual planning state]
\label{def:continual_state}
The continual planning state at time $t$ is the tuple
\[
\Xi_t \triangleq
\big(\mathcal{P}_t,\Omega_t^{\mathrm{fin}},\Omega_t^{\mathrm{run}},
\mathbf{w}_t^{\mathrm{exe}}\big),
\]
which summarizes the currently maintained task structure, the completed
subtasks, the non-preemptive subtasks already in execution, and the realized
high-level behavior.
\end{definition}

Definition~\ref{def:continual_state} is the minimal information required for
online adaptation. In particular, $\mathbf{w}_t^{\mathrm{exe}}$ records the
executed prefix and therefore distinguishes the already discharged temporal
requirements from the requirements that remain to be satisfied. This distinction
is essential in continual settings, since a newly released task should only
interact with the unfinished portion of the current plan.

\subsection{Online product of R-posets}
\label{subsec:online_poset_product}

Assume that a new continual task is released at time $t$ and is represented by
the sc-LTL formula $\varphi_t^{+}$. Following the same automaton-to-poset
abstraction used earlier in the chapter, let
\[
\mathcal{P}_t^{+} \triangleq
(\Omega_t^{+},\preceq_t^{+},\neq_t^{+})
\]
be an R-poset associated with $\varphi_t^{+}$. The goal is to combine
$\mathcal{P}_t^{+}$ with the current global poset $\mathcal{P}_t$ without
reconstructing the complete task model from scratch.

\begin{definition}[Online product of R-posets]
\label{def:online_poset_product}
Given the continual planning state $\Xi_t$ and the incoming R-poset
$\mathcal{P}_t^{+}$, the online product is the set
\begin{equation}
\mathfrak{P}_t^{+} \triangleq
\mathcal{P}_t \otimes_{\mathbf{w}_t^{\mathrm{exe}}} \mathcal{P}_t^{+},
\label{eq:online_poset_product}
\end{equation}
where each element of $\mathfrak{P}_t^{+}$ is an updated R-poset whose
language is compatible with both the unfinished requirements of
$\mathcal{P}_t$ after the executed prefix $\mathbf{w}_t^{\mathrm{exe}}$ and
the newly released task $\mathcal{P}_t^{+}$.
\end{definition}

The product in \eqref{eq:online_poset_product} is computed over the unfinished
portion of $\mathcal{P}_t$. Consequently, the finished subtasks in
$\Omega_t^{\mathrm{fin}}$ are not reconsidered as decision variables. This is
precisely the mechanism that prevents unnecessary recomputation. In contrast
with direct automaton synchronization, the update is performed only at the
level of the remaining subtasks and their partial relations
\cite{liu2024fast}.

The online product contains two conceptual stages. The first stage constructs a
composed subtask set. Each incoming subtask in $\Omega_t^{+}$ is either merged
with a compatible unfinished subtask in $\Omega_t^{\mathrm{rem}}$ or appended
as a new subtask if no compatible image exists. Compatibility is determined by
the execution label and the self-loop condition inherited from the underlying
automata. Therefore, the composition step does not merely concatenate subtasks;
it also identifies cases in which one future execution can satisfy multiple
specifications simultaneously.

To formalize this point, let
$\mu_t : \Omega_t^{+} \rightarrow \bar{\Omega}_t$ be a composition map, where
$\bar{\Omega}_t$ denotes the composed subtask set produced by the current
update. For each $\omega^{+} \in \Omega_t^{+}$, the image $\mu_t(\omega^{+})$
is either an unfinished subtask in $\Omega_t^{\mathrm{rem}}$ with compatible
labels or a newly created subtask appended to $\bar{\Omega}_t$. The mapping
$\mu_t$ thus records how the new temporal specification is embedded into the
updated global poset.

The second stage refines the relations among subtasks. The inherited precedence
and conflict relations from $\mathcal{P}_t$ and $\mathcal{P}_t^{+}$ are first
lifted to the composed subtask set through $\mu_t$. Additional relations are
then introduced when they are required to preserve self-loop guards or to
remove newly induced conflicts. The updated relation sets are therefore written
as
\begin{equation}
\bar{\preceq}_t \triangleq
\preceq_t \cup \mu_t(\preceq_t^{+}) \cup \preceq_t^{\mathrm{add}},
\qquad
\bar{\neq}_t \triangleq
\neq_t \cup \mu_t(\neq_t^{+}) \cup \neq_t^{\mathrm{add}},
\label{eq:online_relation_update}
\end{equation}
where $\preceq_t^{\mathrm{add}}$ and $\neq_t^{\mathrm{add}}$ denote the
extra precedence and conflict relations introduced during refinement.

Equation~\eqref{eq:online_relation_update} is the key structural update in the
continual planning layer. The first two terms preserve the logic already
encoded in the original posets, while the third term resolves incompatibilities
that appear only after composition. As a result, the updated poset
\[
\bar{\mathcal{P}}_t \triangleq
(\bar{\Omega}_t,\bar{\preceq}_t,\bar{\neq}_t)
\]
remains a valid high-level abstraction of the augmented task specification.
Following \eqref{eq:online_poset_product}, one admissible element of
$\mathfrak{P}_t^{+}$ can be selected as the new global poset
$\mathcal{P}_{t^{+}}$. In practice, the first admissible product is often
sufficient for online reaction, whereas a larger portion of
$\mathfrak{P}_t^{+}$ may be explored when additional computation time is
available \cite{liu2024fast}.

\subsection{Receding-horizon planning over the updated poset}
\label{subsec:online_rh_planning}

Once the global poset has been updated, the next step is to adapt the robot
plans. Replanning over the whole remaining subtask set is unnecessary at every
update and would negate the computational advantage gained in
Subsection~\ref{subsec:online_poset_product}. For this reason, the
receding-horizon planning module introduced in the previous section is used
here as the second layer of online adaptation.

Let $\mathcal{P}_t = (\Omega_t,\preceq_t,\neq_t)$ now denote the updated
global poset after the most recent online product. The set of currently ready
subtasks is defined by
\[
\Omega_t^{\mathrm{ready}} \triangleq
\{\omega \in \Omega_t^{\mathrm{rem}} \setminus \Omega_t^{\mathrm{run}}
\mid \forall \omega' \in \Omega_t,\;
\omega' \preceq_t \omega \Rightarrow
\omega' \in \Omega_t^{\mathrm{fin}}\}.
\]
Hence, a subtask is ready only if all its required predecessors have already
been completed and the subtask is not being executed by another robot team.
This definition excludes temporal violations at the high level and also avoids
duplicate dispatch of a currently running subtask.

Given a horizon length $H \in \mathbb{N}_{>0}$, the planner does not optimize
over all remaining subtasks. Instead, it constructs a finite horizon set
\begin{equation}
\Omega_t^{H} \triangleq
\operatorname{Close}_{H}(\Omega_t^{\mathrm{ready}},\preceq_t),
\label{eq:horizon_set}
\end{equation}
where $\operatorname{Close}_{H}$ returns the first $H$ subtasks obtained by
causal expansion from the ready frontier under the precedence relation
$\preceq_t$.

Equation~\eqref{eq:horizon_set} defines the set of subtasks that are exposed to
the receding-horizon optimizer at the current planning cycle. The subtasks in
$\Omega_t^{H}$ may include both immediately executable subtasks and a small
number of near-future subtasks that are needed to preserve local temporal
consistency. In this way, the planner reasons beyond the current frontier while
avoiding the burden of full-horizon optimization.

The non-preemption requirement is imposed through the currently running set
$\Omega_t^{\mathrm{run}}$. More precisely, let $\mathbb{J}_t^{H}$ denote the
set of feasible horizon plans. This set is defined by
\begin{equation}
\mathbb{J}_t^{H} \triangleq
\left\{
\mathbf{J}\ \middle|\
\begin{array}{l}
\mathbf{J}\ \text{preserves all assignments in } \Omega_t^{\mathrm{run}},\\
\mathbf{J}\ \text{satisfies the relations } \preceq_t \text{ and } \neq_t\\
\text{on the subtasks in } \Omega_t^{H},\\
\mathbf{J}\ \text{satisfies the motion, capability, and}\\
\text{collaborative constraints from the previous section}
\end{array}
\right\},
\label{eq:horizon_feasible_set}
\end{equation}
where $\mathbf{J}$ collects the local plans of all robots over the current
planning horizon.

Equation~\eqref{eq:horizon_feasible_set} shows that the present section does
not alter the underlying receding-horizon solver. Instead, it changes the
high-level input to that solver by updating the poset online and by freezing
the subtasks already in execution. Therefore, the local optimization model from
the previous section can be used without modification, except that it is now
applied to the horizon set $\Omega_t^{H}$ derived from the updated global
poset. After the horizon problem is solved, only the first portion of the
resulting plan is committed for execution, and the rest remains revisable in
later planning cycles.

\subsection{Triggering conditions and state handover}
\label{subsec:trigger_handover}

The preceding subsection describes how replanning is performed after an update.
It remains to specify when replanning is triggered and how information is
passed from one planning cycle to the next. Since continual tasks may arrive at
arbitrary times, the planning process is organized around a sequence of
replanning instants $\tau_0,\tau_1,\cdots$.

Three classes of events are monitored online. The first class is a progress
event. This event is generated when a sufficient fraction of the currently
committed horizon subtasks has been completed. The second class is an arrival
event. This event is generated when a new continual task is released or when an
existing high-level request is modified. The third class is an infeasibility
event. This event is generated when the current horizon plan can no longer be
executed because of execution failure, communication loss, or a newly observed
resource conflict.

Let
\[
r_t \triangleq
\frac{
\left| \Omega_t^{\mathrm{fin}} \cap \Omega_{\tau_\ell}^{H} \right|
}{
\left| \Omega_{\tau_\ell}^{H} \right|
}
\]
be the completion ratio of the horizon plan dispatched at the previous
replanning instant $\tau_\ell$, and let $r_{\mathrm{trig}} \in (0,1]$ be the
prescribed progress threshold. The next replanning instant is then defined as
\begin{equation}
\tau_{\ell+1} \triangleq
\inf \left\{
t > \tau_\ell \,\middle|\,
r_t \geq r_{\mathrm{trig}}
\;\vee\;
\mathsf{E}_t^{\mathrm{new}}
\;\vee\;
\mathsf{E}_t^{\mathrm{fail}}
\right\},
\label{eq:replanning_trigger}
\end{equation}
where $\mathsf{E}_t^{\mathrm{new}}$ denotes the arrival of a new or modified
task, and $\mathsf{E}_t^{\mathrm{fail}}$ denotes infeasibility of the current
horizon plan.

Equation~\eqref{eq:replanning_trigger} unifies periodic progress-driven
replanning with event-driven adaptation. The threshold
$r_{\mathrm{trig}} = 1/2$ is often a reasonable choice in practice, since it
balances reactivity against excessive replanning frequency. More importantly,
the arrival and infeasibility events can trigger replanning immediately, which
is necessary in genuinely dynamic environments.

At each replanning instant, the continual planning state is rolled forward using
the observed execution status of the robots and objects. For each robot
$n \in \mathcal{N}$, let $x_n(\tau_{\ell+1})$ be its measured state at the new
planning instant, and let $J_n^{\mathrm{run}}(\tau_{\ell+1})$ be the frozen
prefix of its local plan that contains the subtask currently under execution,
if any. Then the next planning cycle is initialized from the state snapshot
\[
\Xi_{\tau_{\ell+1}} =
\big(\mathcal{P}_{\tau_{\ell+1}},
\Omega_{\tau_{\ell+1}}^{\mathrm{fin}},
\Omega_{\tau_{\ell+1}}^{\mathrm{run}},
\mathbf{w}_{\tau_{\ell+1}}^{\mathrm{exe}}\big),
\]
together with the measured robot and object states. This state handover keeps
the executed portion of the plan fixed, carries over all non-preemptive
subtasks, and makes the remaining assignments available for revision.

The overall continual adaptation loop can now be summarized as follows.
Whenever \eqref{eq:replanning_trigger} is activated, the incoming task formulas
are first absorbed into the global temporal structure through the online poset
product in \eqref{eq:online_poset_product} and the relation update in
\eqref{eq:online_relation_update}. Next, the ready frontier and finite horizon
set are computed according to \eqref{eq:horizon_set}. Finally, the same
receding-horizon planning module as in the previous section is solved over the
feasible set in \eqref{eq:horizon_feasible_set}, subject to the frozen running
subtasks. In this manner, the robots can incorporate newly released continual
tasks while preserving valid progress and avoiding unnecessary recomputation.

%==============================
%==============================
%==============================
\section{Case Study}
\label{sec:ch04_case_study}

The top-down coordination framework is illustrated next on a
representative multi-robot mission with global temporal structure,
hierarchical team organization, and online task updates. The purpose of
the case study is to show how decomposition, assignment, team
formation, and replanning interact in a single execution flow, and to
provide empirical context for the complexity discussion of the
preceding sections.

Consider a mission specified initially at the team level through a global
temporal formula that combines ordered service requirements, allocation
constraints, and limited concurrency. In the first stage, the accepted behavior
of the task automaton is pruned and converted into subtask occurrences together
with a compact precedence structure. This relaxed partially ordered view
separates the logical dependencies that must be preserved from the robot-level
details that can be optimized later. The resulting structure is then used to
assign subtasks to heterogeneous robots while minimizing the mission-level
completion metric. At this stage, the role of top-down synthesis is explicit:
the global specification is not approximated by ad hoc task splitting, but is
systematically refined into allocable components whose temporal relations remain
traceable to the original formula.

In the second stage, the assigned subtasks are grouped into teams and refined
through the hierarchical fleet-coordination machinery of
Section~\ref{ch04:subsec:task}. High-level decomposition determines which robots
should cooperate, team formation enforces capability and redundancy conditions,
and local coordination within each team resolves the motion-level realization of
the assigned subtask sequence. The hierarchy becomes especially valuable when
the fleet is large or heterogeneous, because it prevents the global mission from
being reintroduced as a single fully coupled planning object at every level.
Instead, each layer exposes only the information required by the next one.

Finally, suppose that a new mission element or operator request arrives during
execution. The current global plan is then updated through the online product of
R-posets, and the receding-horizon mechanism of
Section~\ref{ch04:sec:online_continual_tasks} revises only the unexecuted suffix
of the plan. Tasks already in progress are preserved whenever possible, while
new obligations are incorporated into the evolving precedence structure. This
case-study view makes clear that the top-down framework is not merely an
offline decomposition procedure. It is a persistent coordination architecture in
which global temporal structure, hierarchical allocation, and online revision
remain coupled throughout execution.

%==============================
\section{Summary}
\label{ch04:sec:summary}

This chapter developed a top-down framework for global temporal-task
coordination in multi-robot systems. Section~\ref{ch04:sec:po_decomp_assignment}
introduced automaton pruning, task decomposition, relaxed partially ordered
sets, and branch-and-bound assignment for a fixed global mission.
Section~\ref{ch04:subsec:task} extended this framework to large heterogeneous
fleets through hierarchical mission decomposition, abstract team assignment,
team formation, and local coordination within teams.
Section~\ref{ch04:sec:online_continual_tasks} then addressed continual tasks by
updating the task structure online through R-poset products and by applying a
receding-horizon replanning scheme. Human requests were incorporated as special
online task updates within the same framework.

